\section{Introduction}
\label{sec:introduction}

\paragraph{From Single-Pass Rendering to Closed-Loop Visual Creation.}
Visual generation has developed from a specialized media-synthesis capability into a general component of artificial intelligence infrastructure. It now supports the production and editing of images, videos, three-dimensional assets, data-grounded graphics, structured documents, interactive interfaces, and related visual artifacts across design, entertainment, education, science, industry, and other application settings. Advances in generative and multimodal models have substantially improved visual fidelity, instruction and reference adherence, text rendering, controllability, editing, and other task-specific capabilities \cite{he2024llms,yang2025textimage,wang2025imageediting,yin2025spatiotemporal}. A single-pass renderer can now often produce a plausible artifact from a concise instruction and a limited set of conditions.

The expansion from isolated synthesis to multi-step visual creation tasks changes the requirements placed on a visual system. A single artifact or sequence may need to preserve precise spatial relations, target identity, temporal order, editable structure, factual or data-grounded content, physical plausibility, stable interaction, and other cross-step requirements. Surveys, benchmarks, and recent system studies continue to report weaknesses along structural, temporal, physical, and other task-specific dimensions even when individual outputs achieve high perceptual quality \cite{wu2026visualgeneration,yin2025spatiotemporal,avg250925229,feng2026newton}. Their practical consequence is cumulative: an early relation error can invalidate a later edit, identity drift can propagate through a video, an incorrect data mapping can undermine a scientific graphic, and an invalid geometric operation can make a visually plausible model unusable, among other downstream effects. Meeting these requirements calls for creation control: constraint persistence, usable intermediate state, diagnosis, repair, and other forms of trajectory management.

Model scaling, prompt refinement, additional conditioning, repeated sampling, and other output-level refinements can improve individual outputs. A visual creation process remains open-loop when runtime evidence from the evolving artifact or task environment does not alter subsequent decisions. Under this condition, the system cannot use intermediate outcomes to determine whether the current plan remains appropriate, which constraints have been satisfied or violated, where a failure has occurred, or when the task should stop. Local errors and changes in task state may therefore persist or propagate across later edits, frames, components, and other dependent outputs.

Multi-step visual creation requires goal-directed control over the creation trajectory. The system must maintain task-relevant state, inspect intermediate artifacts and execution outcomes, and use the resulting evidence to determine subsequent actions. Recent systems for compositional image generation, image editing, video production, and related visual-creation tasks have begun to implement parts of this process through artifact observation, constraint verification, tool selection, iterative correction, and related operations \cite{wang2024divide,he2024kubrick,gupta2025costa}. These developments illustrate a transition from open-loop sampling toward closed-loop visual creation, in which observations of the evolving artifact and task environment inform later action selection.

In this survey, \emph{agentic visual generation} refers to goal-driven visual creation that operates through a closed loop. Such a system maintains accessible, task-relevant state for an evolving visual artifact, its task environment, and interaction history, then uses runtime observations to select subsequent creation actions. These observations may arise from intermediate artifacts, execution outcomes, external information, user feedback, and other task-relevant sources. The defining behavioral criterion is a demonstrable observation--action dependency: runtime information changes later action selection. Section~\ref{sec:foundations} provides the full operational definition and formalization, establishes the boundary of AVG, and introduces the L1--L5 autonomy levels; Section~\ref{sec:components} examines the analytical components used to organize this control.

\paragraph{Review Motivation and Questions.}
Existing related surveys examine adjacent questions from different starting points. Visual-generation surveys organize the literature around model families, conditioning and editing methods, or particular visual tasks and modalities \cite{yang2025textimage,wang2025imageediting,yin2025spatiotemporal}. The survey of He et al. studies how language models participate in multimodal generation and editing, including through tool-mediated interaction \cite{he2024llms}. Agent surveys instead develop concepts, capabilities, and system designs for language agents or multimodal agents across broad task settings \cite{durante2024agentai,xie2024large,jiang2024rationality,schneider2025generative}. At a different level of abstraction, the recent visual-generation roadmap situates agentic generation within a longer progression toward visual intelligence and world modeling \cite{wu2026visualgeneration}.

These reviews provide essential foundations, but their primary questions differ from the one considered here. To the best of our knowledge, no existing survey is devoted specifically to agentic visual generation as a class of systems in which runtime observations demonstrably influence subsequent visual-creation decisions. This survey uses that behavioral property to organize the literature, then compares how the resulting control mechanisms and supporting evidence vary across visual-creation domains.

The survey addresses six research questions:

\begin{enumerate}
  \item \textbf{Conceptual boundary:} What constitutes agentic visual generation, and how does it relate to single-pass renderers, workflows, and tool-supported visual-creation systems?
  \item \textbf{Analytical components:} Which components organize the connection between runtime observations and subsequent visual-creation decisions?
  \item \textbf{Cross-domain variation:} How are these mechanisms realized across visual-creation domains?
  \item \textbf{Evaluation:} How should the outcomes and process control of AVG systems be evaluated?
  \item \textbf{Training and improvement:} How are AVG systems trained and improved through trainable components, visual trajectories, feedback, experience reuse, and cross-task transfer?
  \item \textbf{Research frontier:} Which open problems and evidence requirements define future AVG research?
\end{enumerate}

\paragraph{Contributions and Organization.}
This survey makes six contributions:

\begin{enumerate}
  \item It defines and formalizes agentic visual generation as closed-loop visual creation and specifies a behavioral boundary based on observation-dependent creation decisions.
  \item It synthesizes a six-component analytical framework comprising goal understanding, specification, and planning; memory; tool; perception; action; and cross-task self-improvement.
  \item It compares how these mechanisms operate across visual products with different representations, constraints, action spaces, sources of verification evidence, and other domain-specific properties.
  \item It analyzes training and improvement in AVG systems from trainable targets and visual trajectory supervision through policy initialization, multimodal and executable-feedback reinforcement learning, experience reuse, skill abstraction, joint optimization, and cross-task transfer.
  \item It organizes evaluation evidence into four complementary levels---artifact quality, goal and constraint satisfaction, trajectory and decision quality, and system and human-centered evaluation---and identifies the remaining gaps in current evaluation practice.
  \item It synthesizes research frontiers for agentic visual generation around expanding the scope of visual agency, increasing the intelligence of visual agents, and establishing reliable evaluation.
\end{enumerate}

The remainder of the survey follows this control-centered organization. Section~\ref{sec:foundations} introduces visual-generation foundations, defines and formalizes AVG, establishes its boundary, and presents the L1--L5 autonomy levels. Section~\ref{sec:components} examines the six analytical components that organize visual-agent systems. Sections~\ref{sec:image-generation}--\ref{sec:cross-domain} compare their realization across visual products and synthesize the domain evidence. Section~\ref{sec:training} examines trajectory supervision, policy initialization, multimodal reinforcement learning, experience reuse, and transfer. Section~\ref{sec:evaluation} organizes evaluation evidence into four complementary levels and identifies remaining evaluation gaps, Section~\ref{sec:frontiers} examines research frontiers, and Section~\ref{sec:conclusion} concludes the survey.
